% GENERATED FILE. Edit canonical lessons, then run npm run generate:books.
% canonical-source-hash: fnv1a64:70d501dcab7fec01
% canonical-lessons: BN-C319-dhvongsaboshesh, BN-C319-murti, BN-C319-khonon-kora, BN-C319-shilalipi, BN-C319-mudra

\chapter{Digging Up the Past}
\label{ch:bn-digging-up-the-past}
\hlchaptermodality{319}

\begin{chapteropening}
\textbf{By the end of this chapter:} \emph{I can talk about ruins, statues, digs and old inscriptions.}

Complete the last lesson of chapter 319: I can talk about ruins, statues, digs and old inscriptions.
\end{chapteropening}

\section[\texorpdfstring{\bn{ধ্বংসাবশেষ} (dhvôngsābôshesh)}{ধ্বংসাবশেষ (dhvôngsābôshesh)}]{\bn{ধ্বংসাবশেষ} (dhvôngsābôshesh) — ruins; remains}

\label{lesson:BN-C319-dhvongsaboshesh}



\begin{quote}
Before the new one: say the Bengali for to found, then the Bengali for liberation.
\end{quote}

\subsection*{You'll want to know: \bn{ধ্বংসাবশেষ}}
\textbf{\bn{ধ্বংসাবশেষ}} — \emph{dhvôngsābôshesh} — \textquotedblleft{}ruins; remains\textquotedblright{}.

\begin{quote}
\bn{নদীর} \bn{পাশে} \bn{একটা} \bn{পুরনো} \bn{শহরের} \bn{ধ্বংসাবশেষ} \bn{আছে।}
\end{quote}

\begin{quote}
\emph{nôdīr pāshe ekṭā purôno shôhôrer dhvôngsābôshesh āchhe.}
\end{quote}

\begin{quote}
Beside the river are the ruins of an old city.
\end{quote}

\begin{cousinweb}
\textbf{\bn{ধ্বংসাবশেষ}} (\emph{dhvôngsābôshesh}) joins \textbf{\bn{ধ্বংস}}, ruin, to Sanskrit \emph{avaśeṣa}, 'what is left' (Monier-Williams), the word in \textbf{\bn{অবশেষে}}: what is left of a ruin.
\end{cousinweb}

\begin{grammarlens}[title={What you've built}]
One more word for telling what happened.
\end{grammarlens}

\subsection*{Your turn}
\begin{itemize}
\raggedright
  \item \emph{Say it:} \emph{dhvôngsābôshesh}
  \item \emph{Say it:} \emph{dhvôngsābôshesh}, once more
  \item \emph{Say it:} \emph{nôdīr pāshe ekṭā purôno shôhôrer dhvôngsābôshesh āchhe.} — Beside the river are the ruins of an old city.
  \item \emph{Recall:} say the Bengali for to found, then the Bengali for liberation, then say \emph{dhvôngsābôshesh} again
\end{itemize}

\begin{tcolorbox}[breakable,colback=teal!4,colframe=teal!35!black,title={Before you move on}]
What does \bn{ধ্বংসাবশেষ} mean? (\textquotedblleft{}Ruins; remains\textquotedblright{}.) Say it once more.
\end{tcolorbox}
\section[\texorpdfstring{\bn{মূর্তি} (mūrti)}{মূর্তি (mūrti)}]{\bn{মূর্তি} (mūrti) — a statue}

\label{lesson:BN-C319-murti}



\begin{quote}
Before the new one: say the Bengali for liberation, then the Bengali for ruins.
\end{quote}

\subsection*{You'll want to know: \bn{মূর্তি}}
\textbf{\bn{মূর্তি}} — \emph{mūrti} — \textquotedblleft{}a statue\textquotedblright{}.

\begin{quote}
\bn{মজুররা} \bn{মাটির} \bn{নিচে} \bn{একটা} \bn{পাথরের} \bn{মূর্তি} \bn{পেল।}
\end{quote}

\begin{quote}
\emph{mojurrā māṭir niche ekṭā pāthôrer mūrti pelô.}
\end{quote}

\begin{quote}
The workmen found a stone statue under the ground.
\end{quote}

\begin{cousinweb}
\textbf{\bn{মূর্তি}} (\emph{mūrti}) is Sanskrit \emph{mūrti}, 'any solid body; an embodiment' (Monier-Williams); a statue gives a god or a hero a solid form.
\end{cousinweb}

\begin{grammarlens}[title={What you've built}]
One more word for telling what happened.
\end{grammarlens}

\subsection*{Your turn}
\begin{itemize}
\raggedright
  \item \emph{Say it:} \emph{mūrti}
  \item \emph{Say it:} \emph{mūrti}, once more
  \item \emph{Say it:} \emph{mojurrā māṭir niche ekṭā pāthôrer mūrti pelô.} — The workmen found a stone statue under the ground.
  \item \emph{Recall:} say the Bengali for liberation, then the Bengali for ruins, then say \emph{mūrti} again
\end{itemize}

\begin{tcolorbox}[breakable,colback=teal!4,colframe=teal!35!black,title={Before you move on}]
What does \bn{মূর্তি} mean? (\textquotedblleft{}A statue\textquotedblright{}.) Say it once more.
\end{tcolorbox}
\section[\texorpdfstring{\bn{খনন} \bn{করা} (khônôn kôrā)}{খনন করা (khônôn kôrā)}]{\bn{খনন} \bn{করা} (khônôn kôrā) — to excavate}

\label{lesson:BN-C319-khonon-kora}



\begin{quote}
Before the new one: say the Bengali for ruins, then the Bengali for a statue.
\end{quote}

\subsection*{You'll want to know: \bn{খনন} \bn{করা}}
\textbf{\bn{খনন} \bn{করা}} — \emph{khônôn kôrā} — \textquotedblleft{}to excavate\textquotedblright{}.

\begin{quote}
\bn{বিজ্ঞানীরা} \bn{দশ} \bn{বছর} \bn{ধরে} \bn{জায়গাটা} \bn{খনন} \bn{করলেন।}
\end{quote}

\begin{quote}
\emph{bigyānīrā dôsh bôchhôr dhôre jāygāṭā khônôn kôrlen.}
\end{quote}

\begin{quote}
The scientists excavated the site for ten years.
\end{quote}

Past (he/she …): \bn{তিনি} \bn{পুরনো} \bn{শহরটা} \bn{খনন} \bn{করলেন}, \emph{tini purôno shôhôrṭā khônôn kôrlen}.

\begin{cousinweb}
\textbf{\bn{খনন} \bn{করা}} (\emph{khônôn kôrā}) puts Sanskrit \emph{khanana}, 'the act of digging or excavating' (Monier-Williams), before \bn{করা}; \textbf{\bn{খোঁড়া}} is the everyday word for digging.
\end{cousinweb}

\begin{grammarlens}[title={What you've built}]
One more word for telling what happened.
\end{grammarlens}

\subsection*{Your turn}
\begin{itemize}
\raggedright
  \item \emph{Say it:} \emph{khônôn kôrā}
  \item \emph{Say it:} \emph{khônôn kôrā}, once more
  \item \emph{Say it:} \emph{bigyānīrā dôsh bôchhôr dhôre jāygāṭā khônôn kôrlen.} — The scientists excavated the site for ten years.
  \item \emph{Recall:} say the Bengali for ruins, then the Bengali for a statue, then say \emph{khônôn kôrā} again
\end{itemize}

\begin{tcolorbox}[breakable,colback=teal!4,colframe=teal!35!black,title={Before you move on}]
What does \bn{খনন} \bn{করা} mean? (\textquotedblleft{}To excavate\textquotedblright{}.) Say it once more.
\end{tcolorbox}
\section[\texorpdfstring{\bn{শিলালিপি} (shilālipi)}{শিলালিপি (shilālipi)}]{\bn{শিলালিপি} (shilālipi) — an inscription (on stone)}

\label{lesson:BN-C319-shilalipi}



\begin{quote}
Before the new one: say the Bengali for a statue, then the Bengali for to excavate.
\end{quote}

\subsection*{You'll want to know: \bn{শিলালিপি}}
\textbf{\bn{শিলালিপি}} — \emph{shilālipi} — \textquotedblleft{}an inscription (on stone)\textquotedblright{}.

\begin{quote}
\bn{শিলালিপিটা} \bn{পড়ে} \bn{বিজ্ঞানীরা} \bn{রাজার} \bn{নাম} \bn{জানলেন।}
\end{quote}

\begin{quote}
\emph{shilālipiṭā pôṛe bigyānīrā rājār nām jānlen.}
\end{quote}

\begin{quote}
By reading the inscription the scientists learned the king's name.
\end{quote}

\begin{cousinweb}
\textbf{\bn{শিলালিপি}} (\emph{shilālipi}) joins Sanskrit \emph{śilā}, 'a stone, a rock', to \emph{lipi}, 'writing, a script' (Monier-Williams): writing cut in stone.
\end{cousinweb}

\begin{grammarlens}[title={What you've built}]
One more word for telling what happened.
\end{grammarlens}

\subsection*{Your turn}
\begin{itemize}
\raggedright
  \item \emph{Say it:} \emph{shilālipi}
  \item \emph{Say it:} \emph{shilālipi}, once more
  \item \emph{Say it:} \emph{shilālipiṭā pôṛe bigyānīrā rājār nām jānlen.} — By reading the inscription the scientists learned the king's name.
  \item \emph{Recall:} say the Bengali for a statue, then the Bengali for to excavate, then say \emph{shilālipi} again
\end{itemize}

\begin{tcolorbox}[breakable,colback=teal!4,colframe=teal!35!black,title={Before you move on}]
What does \bn{শিলালিপি} mean? (\textquotedblleft{}An inscription (on stone)\textquotedblright{}.) Say it once more.
\end{tcolorbox}
\section[\texorpdfstring{\bn{মুদ্রা} (mudrā)}{মুদ্রা (mudrā)}]{\bn{মুদ্রা} (mudrā) — currency; a coin (of old)}

\label{lesson:BN-C319-mudra}



\begin{quote}
Before the new one: say the Bengali for to excavate, then the Bengali for an inscription.
\end{quote}

\subsection*{You'll want to know: \bn{মুদ্রা}}
\textbf{\bn{মুদ্রা}} — \emph{mudrā} — \textquotedblleft{}currency; a coin (of old)\textquotedblright{}.

\begin{quote}
\bn{এই} \bn{প্রাচীন} \bn{মুদ্রাগুলো} \bn{দুই} \bn{হাজার} \bn{বছরের} \bn{পুরনো।}
\end{quote}

\begin{quote}
\emph{ei prāchīn mudrāgulo dui hājār bôchhôrer purôno.}
\end{quote}

\begin{quote}
These ancient coins are two thousand years old.
\end{quote}

\begin{cousinweb}
\textbf{\bn{মুদ্রা}} (\emph{mudrā}) is Sanskrit \emph{mudrā}, a seal or a stamp, and so a stamped coin (Monier-Williams); today it also means a country's currency.
\end{cousinweb}

\begin{grammarlens}[title={What you've built}]
One more word for telling what happened.
\end{grammarlens}

\subsection*{Your turn}
\begin{itemize}
\raggedright
  \item \emph{Say it:} \emph{mudrā}
  \item \emph{Say it:} \emph{mudrā}, once more
  \item \emph{Say it:} \emph{ei prāchīn mudrāgulo dui hājār bôchhôrer purôno.} — These ancient coins are two thousand years old.
  \item \emph{Recall:} say the Bengali for to excavate, then the Bengali for an inscription, then say \emph{mudrā} again
\end{itemize}

\begin{tcolorbox}[breakable,colback=teal!4,colframe=teal!35!black,title={Before you move on}]
What does \bn{মুদ্রা} mean? (\textquotedblleft{}Currency; a coin (of old)\textquotedblright{}.) Say it once more.
\end{tcolorbox}
